\subsection{相伴分次模中元素族的提升}

设 A 为滤过交换环，\((\mathbf{A}_{n})_{n\in\mathbf{Z}}\) 为其滤过，C 为 \(A_{0}\) 的子环，满足 \(C\cap A_{1}=\{0\}\)。于是，\(\mathbf{A}_{0}\to\mathbf{A}_{0}/\mathbf{A}_{1}=\mathrm{gr}_{0}(\mathbf{A})\) 在 C 上的限制是单射，这使我们能够把 C 识别为 \(\mathrm{gr}_{0}(\mathbf{A})\) 的子模；类似情况下我们通常都这样做。若 \(A_{1}\neq A_{0}\)，且 K 是 \(A_{0}\) 的任意子域，则 \(K\cap A_{1}=\{0\}\)，因为 \(K\cap A_{1}\) 是 K 的不含 1 的理想；于是 K 可识别为 \(\mathrm{gr}_{0}(\mathbf{A})\) 的子域。

	命题 10. 设 A 为滤过交换环，\((\mathbf{A}_n)\) 为其滤过；假设存在子环 \(\mathbf{C}\)，它包含于 \(\mathbf{A}_0\)，使得 \(\mathbf{C} \cap \mathbf{A}_1 = \{0\}\)，并且将 \(\mathbf{C}\) 识别为 \(\operatorname{gr}_0(\mathbf{A})\) 的子环。设 \((x_i)_{1 \leqslant i \leqslant q}\) 是 \(\mathbf{A}\) 中的有限元素族；假设 \(x_i \in \mathbf{A}_{n_i}\) 对 \(1 \leqslant i \leqslant q\) 成立，并令 \(\xi_i\) 为 \(x_i\) 在 \(\operatorname{gr}_{n_i}(\mathbf{A})\) 中的类，对 \(1 \leqslant i \leqslant q\) 成立。

\begin{enumerate}
\def\labelenumi{(\roman{enumi})}
\item
	  若元素族 \((\xi_{i})\) 是 \(\operatorname{gr}(\mathbf{A})\) 中关于 \(\mathbf{C}\) 代数自由的元素族，则元素族 \((x_{i})\) 在 \(\mathbf{C}\) 上代数自由。
\item
  若 \(\mathbf{A}\) 上的滤过穷竭且离散，且 \((\xi_{i})\) 是 C-代数 \(\operatorname{gr}(\mathbf{A})\) 的生成系，则 \((x_{i})\) 是 C-代数 \(\mathbf{A}\) 的生成系。
\end{enumerate}

	令 \(\mathbf{A}'\) 为 \(\mathbf{C}[\mathbf{X}_1, \ldots, \mathbf{X}_q]\) 上的 \(\mathbf{C}\)-多项式代数；给 \(\mathbf{A}'\) 赋予分次 \((\mathbf{A}_n')\)，其类型为 \(\mathbf{Z}\)，其中 \(\mathbf{A}_n'\) 是单项式 \(\mathbf{X}_1^{s(1)} \ldots \mathbf{X}_q^{s(q)}\) 的 C-线性组合所成的集合，满足 \(\sum_{i=1}^{q} n_i s(i) = n\)。令 \(u\) 为同态 \(f \mapsto f(x_1, \ldots, x_q)\)，从 C-代数 \(\mathbf{A}'\) 到 C-代数 \(\mathbf{A}\)；按定义，\(u(\mathbf{A}_n') \subset \mathbf{A}_n\) 对所有 \(n \in \mathbf{Z}\) 成立，因而 \(u\) 与滤过相容（给 \(\mathbf{A}'\) 赋予与其分次环结构相伴的滤过环结构，参见~第 1 节，例 1）。因此，(i) 的假设意味着

\[
\operatorname{gr} (u): \mathrm{A} ^ {\prime} = \operatorname{gr} (\mathrm{A} ^ {\prime}) \rightarrow \operatorname{gr} (\mathrm{A})
\]

是单射；由于 A' 上的滤过穷竭且分离，可以应用第 8 节定理 1 的推论 1，故 u 是单射，这证明了 (i) 的结论。同样，关于 \((\xi_{i})\) 的假设 (ii) 意味着 gr(u) 是满射；由于 A 是离散的且其滤过穷竭，可以应用第 8 节定理 1 的推论 2，故 u 是满射，这证明了 (ii) 的结论。

	命题 11. 设 A 为完备 Hausdorff 滤过交换环，C 为 \(\mathbf{A}_0\) 的子环，满足 \(\mathbf{C} \cap \mathbf{A}_1 = 0\)，且 \((x_i)_{1 \leqslant i \leqslant q}\) 为 A 中的有限元素族，满足 \(x_i \in \mathbf{A}_{n_i}\)，其中 \(n_i > 0\) 对 \(1 \leqslant i \leqslant q\) 成立；令 \(\xi_i\) 为 \(x_i\) 在 \(\operatorname{gr}_{n_i}(\mathbf{A})\) 中的类，对 \(1 \leqslant i \leqslant q\) 成立。(i) 存在唯一的 C-同态 \(v\)，从形式幂级数代数 \(\mathbf{A}'' = \mathbf{C}[[\mathbf{X}_1, \ldots, \mathbf{X}_q]]\) 映到 A，并满足 \(v(\mathbf{X}_i) = x_i\) 对 \(1 \leqslant i \leqslant q\) 成立。

\begin{enumerate}
\def\labelenumi{(\roman{enumi})}
\setcounter{enumi}{1}
\item
  若族 \((\xi_{i})\) 在 \(\mathbf{C}\) 上代数自由，则同态 \(v\) 是单射。
\item
  若 \(\mathbf{A}\) 上的滤过穷竭，且族 \((\xi_{i})\) 是 C-代数 \(\operatorname{gr}(\mathbf{A})\) 的生成系，则同态 \(v\) 是满射。
\end{enumerate}

	由于 \(n_i \geqslant 1\) 对所有 \(i, \sum_{i=1}^{q} n_i s(i) \geqslant \sum_{i=1}^{q} s(i)\) 的情形，对每个单项式 \(X_1^{s(1)} \ldots X_q^{s(q)}\) 都有上述关系；另一方面，不等式 \(\sum_{i=1}^{q} n_i s(i) \leqslant r. \sum_{i=1}^{q} s(i)\) 在 \(r\) 为各 \(n_i\) 中最大的一个时成立。若 \(A_n''\) 表示由这样的形式幂级数所成的集合：其非零项 \(a_s X_1^{s(1)} \ldots X_q^{s(q)}\) 满足 \(\sum_{i=1}^{q} n_i s(i) \geqslant n\)，则由第 6 节命题 6 的推论可知，\(A''\) 在穷竭滤过 \((A_n'')\) 下是 Hausdorff 且完备的，并且

\[
\mathrm{A} ^ {\prime} = \mathrm{C} [ \mathrm{X} _ {1}, \dots , \mathrm{X} _ {q} ]
\]

在 A'' 中稠密；此外，证明命题 10 时定义的同态 u 在 A' 上连续，并且由于 A 是 Hausdorff 且完备的，可以唯一延拓为连续同态 v: A''→A（《一般拓扑学》，第 III 章，§ 3，第 3 节，命题 5），这证明了 (i)；又有 gr(A'') = gr(A') 且 gr(v) = gr(u)；根据 A 的假设，(ii) 和 (iii) 分别由第 8 节定理 1 的推论 1 和 2 得出。

命题的 (ii)（分别 (iii)）的结论有时表述为：族 \((x_{i})\) 在 C 上形式自由（分别是 A 的形式生成系）。

命题 12. 设 A 为滤过环，E 为滤过 A-模，\((\mathbf{A}_{n})\) 和 \((\mathbf{E}_{n})\) 分别是 A 和 E 上的滤过。假设 A 完备且滤过 \((\mathbf{E}_{n})\) 穷竭并分离。设 \((x_{i})_{i\in\mathbf{I}}\) 是 E 中的有限元素族，并对 \(i\in I\) 令 \(n(i)\) 为满足 \(x_{i}\in\mathbf{E}_{n(i)}\) 的整数；最后令 \(\xi_{i}\) 为 \(x_{i}\) 在 \(\mathrm{gr}_{n(i)}(\mathbf{E})\) 中的类。则若 \((\xi_{i})\) 是 gr(A)-模 gr(E) 的生成系，\((x_{i})\) 是 A-模 E 的生成系。

	在 A-模 \(\mathbf{L} = \mathbf{A}_s^{\mathbf{I}}\) 中，令 \(\mathbf{L}_n\) 表示由 \((a_i)\) 所成的集合，其中 \(a_i \in \mathbf{A}_{n - n(i)}\) 对所有 \(i \in \mathbf{I}\) 成立；若 \(p\) 和 \(q\) 分别是 \(n(i)\) 中最小和最大的，则 \(\mathbf{A}_{n - q}^{\mathbf{I}} \supset \mathbf{L}_n \supset \mathbf{A}_{n - p}^{\mathbf{I}}\)，并且由 \(\mathbf{L}\) 上的定义 \((\mathbf{L}_n)\) 所确定的拓扑就是乘积拓扑；因此 \(\mathbf{L}\) 是完备滤过 A-模。由于 \(\mathbf{L}\) 是自由模，存在一个

A-线性映射 \(u: \mathbf{L} \to \mathbf{E}\)，满足 \(u((a_i)) = \sum_{i \in \mathbf{I}} a_i x_i\)，且它显然与滤过相容；我们必须证明 \(u\) 是满射，而根据第 8 节定理 1 的推论 2，为此只需证明

\[
\operatorname{gr} (u): \operatorname{gr} (\mathbf {L}) \rightarrow \operatorname{gr} (\mathbf {E})
\]

	是满射；或者等价地，对所有 \(x \in \mathbf{E}_n\)，存在一个族 \((a_i)\)，使 \(a_i \in \mathbf{A}_{n-n(i)}\) 对所有 \(i \in I\) 成立，且 \(x \equiv \sum_{i \in I} a_i x_i \pmod{\mathbf{E}_{n+1}}\)。令 \(\xi\) 为 \(x\) 在 \(\mathrm{gr}_n(\mathbf{E})\) 中的类；由于 \(\xi_i\) 生成 \(\mathrm{gr}(\mathbf{A})\)-模 \(\mathrm{gr}(\mathbf{E})\)，存在 \(\alpha_i \in \mathrm{gr}(\mathbf{A})\)，使得 \(\xi = \sum_{i \in I} \alpha_i \xi_i\)，并且假设 \(\alpha_i \in \mathrm{gr}_{n-n(i)}(\mathbf{A})\)；必要时可将 \(\alpha_i\) 替换为其 \(n - n(i)\) 次齐次分量。于是 \(\alpha_i\) 是某个 \(a_i \in \mathbf{A}_{n-n(i)}\) 的像，族 \((a_i)\) 具有所要求的性质。

推论 1. 设 A 为完备滤过环，E 为滤过 A-模，且其滤过穷竭并分离。若 gr(E) 是有限生成的（分别为 Noether 的）gr(A)-模，则 E 是有限生成的（分别为 Noether 的）A-模。

若 \(\mathrm{gr}(\mathrm{E})\) 有限生成，则存在有限齐次生成系，由命题 12 可知 E 有限生成。现在假设 \(\mathrm{gr}(\mathrm{E})\) 是 Noether 的，并设 F 为 E 的子模；由 E 上的滤过在 F 上诱导的滤过穷竭且分离，而 \(\mathrm{gr}(\mathrm{F})\) 可识别为 gr(E) 的一个 gr(A)-子模（第 4 节，命题 2），故由假设它有限生成；于是 F 是有限生成 A-模，从而 E 是 Noether 模。

推论 2. 设 A 为完备 Hausdorff 滤过环，且其滤过穷竭。若 gr(A) 是左 Noether 环，则 A 也是左 Noether 环。

只需将推论 1 应用于 \(\mathbf{E} = \mathbf{A}_s\) 即可。

推论 3. 设 A 为完备滤过环，\((\mathbf{A}_{n})\) 为其滤过，E 为 Hausdorff 滤过 A-模，\((\mathbf{E}_{n})\) 为其滤过，F 为 E 的有限生成子模；假设 \(A_{0} = A\) 且 \(E_{0} = E\)。

\begin{enumerate}
\def\labelenumi{(\roman{enumi})}
\item
  若对所有 \(k \geqslant 0\) 都有 \(\mathbf{E}_k = \mathbf{E}_{k+1} + \mathbf{A}_k\mathbf{F}\)，则 \(\mathbf{F} = \mathbf{E}\)。
\item
  若进一步假设 \(\mathbf{E}\) 上的滤过由 \(\mathbf{A}\) 上的滤过导出（第 1 节，例 2），则关系 \(\mathbf{E} = \mathbf{E}_1 + \mathbf{F}\) 蕴含 \(\mathbf{F} = \mathbf{E}\)。
\end{enumerate}

	令 \(\xi_{i}(1\leqslant i\leqslant n)\) 为 \(\mathbf{E}_1\) 模下 \(\mathbf{F}\) 的一个有限生成系的类。由给定假设，对所有 \(k\geqslant 0\)，\(\mathrm{gr}_k(\mathrm{E})\) 的每个元素

	都可写成 \(\sum_{i=1}^{\infty} \alpha_i \xi_i\) 的形式，其中 \(\alpha_i \in \mathrm{gr}(A)\)；因此 \(\xi_i\) 生成 \(\mathrm{gr}(A)\)-模 \(\mathrm{gr}(E)\)，由命题 12 得到 (i)。若 \(E\) 上的滤过由 \(A\) 上的滤过导出，则关系 \(E = E_1 + F\) 蕴含

\[
\begin{array}{r} \mathbf {E} _ {k} = \mathbf {A} _ {k} \mathbf {E} = \mathbf {A} _ {k} \mathbf {E} _ {1} + \mathbf {A} _ {k} \mathbf {F} = \mathbf {A} _ {k} \mathbf {A} _ {1} \mathbf {E} + \mathbf {A} _ {k} \mathbf {F} \subset \mathbf {A} _ {k + 1} \mathbf {E} + \mathbf {A} _ {k} \mathbf {F} \\ = \mathbf {E} _ {k + 1} + \mathbf {A} _ {k} \mathbf {F} \subset \mathbf {E} _ {k}, \end{array}
\]

从而得到 (ii)。

命题 13. 设 A 为环，m 为包含在 A 的 Jacobson 根中的 A 的双边理想，E 为 A-模。给 A 和 E 赋予 m-进滤过（第 1 节，例 3）。假设下列条件之一成立：

\begin{enumerate}
\def\labelenumi{(\alph{enumi})}
\item
  E 是有限生成 A-模且 A 是 Hausdorff 的；
\item
  m 是幂零的。
\end{enumerate}

E 是自由 A-模的充要条件是：E/mE 是自由 (A/m)-模，且 E 满足性质 (GR)（第 8 节）。

	若 E 是自由 A-模，且 \((e_{\lambda})\) 是 E 的一组基，则 \(m^{k}E\) 是 E 中各子模 \(m^{k}e_{\lambda}\) 的直和，对所有 \(k \geqslant 0\) 成立（《代数》，第 II 章，§3，第 7 节，注）；于是 \(m^{k}E/m^{k+1}E\) 可识别为各 \(m^{k}e_{\lambda}/m^{k+1}e_{\lambda}\) 的直和（《代数》，第 II 章，§1，第 6 节，命题 7）。我们首先（当 k = 0 时）得到，各 \(1 \otimes e_{\lambda}\) 是各 \(e_{\lambda}\) 的类，并且它们位于 E/mE = (A/m) \(\otimes_{A} E\) 中，构成 (A/m)-模 E/mE 的一组基，因为典范映射

\[
\left(\mathfrak {m} ^ {k} / \mathfrak {m} ^ {k + 1}\right) \otimes_ {\mathbf {A}} (\mathrm{E} / \mathfrak {m} \mathrm{E}) \rightarrow \mathfrak {m} ^ {k} \mathrm{E} / \mathfrak {m} ^ {k + 1} \mathrm{E}
\]

	对于所有 \(k \geqslant 0\) 都是双射；因此 \(\gamma_{\mathbf{E}}\) 是双射。注意，证明的这一部分既未用到条件 (a)，也未用到条件 (b)。

	反之，设命题条件成立，并令 \((x_{\iota})_{\iota \in I}\) 是 E 中一族元素，其模 mE 的类构成 (A/m)-模 E/mE 的一组基；令 L 为自由 A-模 \(\mathrm{A}_s^{(I)}\)，\((f_{\iota})_{\iota \in I}\) 为其典范基，并令 \(u: \mathbf{L} \to \mathbf{E}\) 为满足 \(u(f_{\iota}) = x_{\iota}\)（对所有 \(\iota \in I\)）的 A-线性映射。命题的假设已蕴含 \(u\) 为满射（第 II 章，§3，第 2 节，命题 4 的推论 1），所以只需证明 \(u\) 为单射。现在，条件 (a) 和 (b) 中的每一个都蕴含 A 是 Hausdorff 的，因而带 m-进滤过的 L 也具有此性质，因为 \(\mathfrak{m}^k\mathbf{L} = (\mathfrak{m}^k)^{(I)}\)（《代数》，第 II 章，§3，第 7 节，注），且 \(\operatorname{gr}(\mathbf{L})\) 可根据证明的第一部分识别为 \(\operatorname{gr}(\mathbf{A}) \otimes_{\mathbf{A} / \mathfrak{m}} (\mathbf{L} / \mathfrak{m}\mathbf{L})\)；同态 \(u\) 与滤过相容，并且可以写成 \(\operatorname{gr}(u) = \gamma_{\mathbb{E}} \circ v\)，其中 \(v\) 是把 \(\operatorname{gr}(\mathbf{L})\) 双射到

\[
\operatorname{gr} (\mathbf {A}) \otimes_ {\mathbf {A} / \mathfrak {m}} (\mathbf {E} / \mathfrak {m} \mathbf {E})
\]

	将 \(f_{t}\) 模 mM 的类映到 \(1 \otimes \bar{x}_{t}\)，其中 \(\bar{x}_{t}\) 是 \(x_{t}\) 模 mE 的类。于是该假设蕴含 \(\operatorname{gr}(u)\) 是单射，结论由第 8 节定理 1 的推论 1 得出。

